% Zone „Das kennst du schon“ – aus bank/skalarprodukt-und-winkel/zone.jsonl (zusammenbau v0.1)
\einheitenkopf[zone][Kennst du schon]{Das kennst du schon}
\setcounter{aufgabe}{0}

%% TODO Titel: Ich-kann-Satz fehlt in der Bank (Platzhalter: Kettenname) – Nr. 1
%% TODO Anweisung: ein Satz über der ersten Teilaufgabe fehlt in der Bank – Nr. 1
\begin{aufgabe}{Skalarprodukt in Koordinaten berechnen, Beträge bilden}
\begin{teile}
\teil $\vec a = (1 \mid 2 \mid 3)$, $\vec b = (3 \mid 1 \mid 1)$. Berechne $\vec a \circ \vec b$. $\vec a \circ \vec b =$ \leerfeld
\teil $\vec a = (-2 \mid 0{,}5 \mid 3)$, $\vec b = (4 \mid -6 \mid 1)$. Berechne $\vec a \circ \vec b$. $\vec a \circ \vec b =$ \leerfeld
\end{teile}
\end{aufgabe}

%% TODO Titel: Ich-kann-Satz fehlt in der Bank (Platzhalter: Kettenname) – Nr. 2
%% TODO Anweisung: ein Satz über der ersten Teilaufgabe fehlt in der Bank – Nr. 2
\begin{aufgabe}{Normalenvektoren aus Koordinatengleichungen ablesen und aus Spannvektoren bestimmen}
\begin{teile}
\teil Gib einen Normalenvektor der Ebene $3x + 2y + 5z = 7$ an. $\vec n = (\,\_\_ \mid \_\_ \mid \_\_\,)$
\teil Eine Ebene wird von $\vec u = (1 \mid 0 \mid 2)$ und $\vec v = (0 \mid 1 \mid 3)$ aufgespannt. Bestimme einen Normalenvektor. $\vec n = (\,\_\_ \mid \_\_ \mid \_\_\,)$
\end{teile}
\end{aufgabe}

%% TODO Titel: Ich-kann-Satz fehlt in der Bank (Platzhalter: Kettenname) – Nr. 3
%% TODO Anweisung: ein Satz über der ersten Teilaufgabe fehlt in der Bank – Nr. 3
\begin{aufgabe}{Richtungsvektoren aus Geraden und Kanten bilden}
\begin{teile}
\teil $A(1 \mid 2 \mid 0)$, $B(4 \mid 3 \mid 5)$. Gib den Vektor $\overrightarrow{AB}$ an. $\overrightarrow{AB} = (\,\_\_ \mid \_\_ \mid \_\_\,)$
\teil $P(-2 \mid 1{,}5 \mid 3)$, $Q(1 \mid -0{,}5 \mid 7)$. Gib den Vektor $\overrightarrow{PQ}$ an. $\overrightarrow{PQ} = (\,\_\_ \mid \_\_ \mid \_\_\,)$
\end{teile}
\end{aufgabe}

%% TODO Titel: Ich-kann-Satz fehlt in der Bank (Platzhalter: Kettenname) – Nr. 4
%% TODO Anweisung: ein Satz über der ersten Teilaufgabe fehlt in der Bank – Nr. 4
\begin{aufgabe}{Kosinus und Tangens am Taschenrechner auswerten und die Umkehrfunktion anwenden (Grad-Modus)}
\begin{teile}
\teil $\mathrm{cos}\,\varphi = 0{,}5$ mit $0^\circ \le \varphi \le 180^\circ$. Bestimme $\varphi$. $\varphi =$ \leerfeld
\teil $\mathrm{cos}\,\varphi = \frac{2}{9}$ mit $0^\circ \le \varphi \le 180^\circ$. Bestimme $\varphi$ auf eine Stelle. $\varphi \approx$ \leerfeld
\end{teile}
\end{aufgabe}

%% TODO Titel: Ich-kann-Satz fehlt in der Bank (Platzhalter: Kettenname) – Nr. 5
%% TODO Anweisung: ein Satz über der ersten Teilaufgabe fehlt in der Bank – Nr. 5
\begin{aufgabe}{Neben-, Komplement- und Innenwinkel unterscheiden und ergänzen}
\begin{teile}
\teil Ein Winkel ist $70^\circ$ groß. Wie groß ist sein Nebenwinkel? \leerfeld
\teil Ein Winkel ist $73{,}4^\circ$ groß. Wie groß ist sein Nebenwinkel? \leerfeld
\end{teile}
\end{aufgabe}

%% TODO Titel: Ich-kann-Satz fehlt in der Bank (Platzhalter: Kettenname) – Nr. 6
%% TODO Anweisung: ein Satz über der ersten Teilaufgabe fehlt in der Bank – Nr. 6
\begin{aufgabe}{Kreisumfang und Anteile, Prozentrechnung}
\begin{teile}
\teil Ein Kreis hat den Umfang $40$ cm. Wie lang ist der Bogen zum Mittelpunktswinkel $90^\circ$? \leerfeld[cm]
\teil Ein Kreis hat den Radius $5$ cm. Wie lang ist der Bogen zum Mittelpunktswinkel $72^\circ$? Runde auf eine Stelle. \leerfeld[cm]
\end{teile}
\end{aufgabe}

%% TODO Titel: Ich-kann-Satz fehlt in der Bank (Platzhalter: Kettenname) – Nr. 7
%% TODO Anweisung: ein Satz über der ersten Teilaufgabe fehlt in der Bank – Nr. 7
\begin{aufgabe}{Neben-, Komplement- und Innenwinkel unterscheiden und ergänzen – Fehler finden}
\begin{teile}
\teil Mia soll den Komplementwinkel von $34^\circ$ angeben und rechnet $180^\circ - 34^\circ = 146^\circ$. Finde den Fehler und rechne richtig.
\end{teile}
\end{aufgabe}

%% TODO Titel: Ich-kann-Satz fehlt in der Bank (Platzhalter: Kettenname) – Nr. 8
%% TODO Anweisung: ein Satz über der ersten Teilaufgabe fehlt in der Bank – Nr. 8
\begin{aufgabe}{Neben-, Komplement- und Innenwinkel unterscheiden und ergänzen – selbst rechnen}
\begin{teile}
\teil Ein Winkel ist $41^\circ$ groß. Wie groß ist sein Komplementwinkel? \leerfeld
\end{teile}
\end{aufgabe}
